\documentclass[12pt,a4paper]{article}
\usepackage[margin=1in]{geometry}
\usepackage{booktabs}
\usepackage{longtable}
\usepackage{enumitem}
\usepackage{listings}
\usepackage{xcolor}
\usepackage{hyperref}

\lstset{
    basicstyle=\ttfamily\small,
    backgroundcolor=\color{gray!10},
    frame=single,
    breaklines=true,
    columns=fullflexible
}

\title{\textbf{Testing Report: Custom System Calls} \\ 
\large Group 4 \\ 
Assignment 2}
\author{
\begin{tabular}{|l|l|l|}
\hline
\textbf{Name} & \textbf{Roll Number} \\
\hline
Payidiparthy Sri Surya Naga Hadwik & 23CS10052 \\
Chintada Yesheeth Sree Narayana & 23CS10013 \\
Dasari Veera Venkata Abhinav Teja & 23CS10017 \\
\hline
\end{tabular} 
}
\date{}

\begin{document}
\maketitle

\section*{Environment}
\begin{itemize}[nosep]
    \item \textbf{OS:} Ubuntu 24.04.4 LTS (running in VirtualBox)
    \item \textbf{Kernel:} Linux 6.1.6 (compiled from source)
    \item \textbf{Compiler:} GCC (Ubuntu default)
    \item \textbf{Architecture:} x86\_64
\end{itemize}

\section*{Task A: Process Information (\texttt{proc\_info})}

\subsection*{Test 1: Current Process Info (PID Mode)}
\textbf{Input:} \texttt{get\_proc\_info(getpid())} \\
\textbf{Expected:} Returns valid process information for the running test process. \\
\textbf{Result:} PASS

\begin{lstlisting}
Parent PID:        3616
State:             0           (TASK_RUNNING)
Priority:          120         (default for SCHED_NORMAL)
Scheduling Class:  SCHED_NORMAL
Child Processes:   0
Sibling Processes: 1
\end{lstlisting}

\textbf{Verification:} State 0 (RUNNING) is correct since the process is actively executing the syscall. Priority 120 is the default for normal processes. SCHED\_NORMAL is the expected scheduler for user programs.

\subsection*{Test 2: Init Process (PID 1)}
\textbf{Input:} \texttt{get\_proc\_info(1)} \\
\textbf{Expected:} Returns information about the init/systemd process. \\
\textbf{Result:} PASS

\begin{lstlisting}
Parent PID:        0           (init's parent is the kernel)
State:             1           (TASK_INTERRUPTIBLE)
Priority:          120
Scheduling Class:  SCHED_NORMAL
Child Processes:   33
Sibling Processes: 2
\end{lstlisting}

\textbf{Verification:} PID 1 (systemd) correctly shows parent PID 0 (the kernel idle process). The high child count reflects the system services spawned by init.

\subsection*{Test 3: System-Wide Info}
\textbf{Input:} \texttt{get\_system\_info()} \\
\textbf{Expected:} Returns aggregate system statistics with consistent process counts. \\
\textbf{Result:} PASS

\begin{lstlisting}
Total Processes:       327
Running:               1
Interruptible:         219
Uninterruptible:       0
RT Class:              57
Fair Class:            270
In CFS Runqueue:       1
Min Vruntime PID:      4436
Min Vruntime:          48932457838
Total CFS Load:        1048576
\end{lstlisting}

\textbf{Verification:} RT (57) + Fair (270) = 327 = Total Processes, confirming counts are consistent. Only 1 process is running (the test program itself), which is expected on a mostly idle system. The min vruntime PID matches the test program's PID.

\subsection*{Test 4: Invalid PID}
\textbf{Input:} \texttt{syscall(451, 99999, PROC\_INFO\_PID, buf, sizeof(buf))} \\
\textbf{Expected:} Returns $-1$ with \texttt{errno} set to \texttt{ESRCH}. \\
\textbf{Result:} PASS --- \texttt{Error: No such process}

\subsection*{Test 5: Invalid Flags}
\textbf{Input:} \texttt{syscall(451, 1, 99, buf, sizeof(buf))} \\
\textbf{Expected:} Returns $-1$ with \texttt{errno} set to \texttt{EINVAL}. \\
\textbf{Result:} PASS --- \texttt{Error: Invalid argument}

\subsection*{Test 6: Buffer Too Small}
\textbf{Input:} \texttt{syscall(451, 1, PROC\_INFO\_PID, buf, 1)} \\
\textbf{Expected:} Returns $-1$ with \texttt{errno} set to \texttt{ENOSPC}. \\
\textbf{Result:} PASS --- \texttt{Error: No space left on device}

\section*{Task B: Context Switch Tracker}

\subsection*{Test 1: Register Current Process}
\textbf{Input:} \texttt{register\_pid(getpid())} \\
\textbf{Expected:} Returns 0, process is added to the monitored list. \\
\textbf{Result:} PASS --- \texttt{Registered PID 3859 successfully}

\subsection*{Test 2: Register PID 1 (systemd)}
\textbf{Input:} \texttt{register\_pid(1)} \\
\textbf{Expected:} Returns 0, PID 1 is added to the monitored list. \\
\textbf{Result:} PASS --- \texttt{Registered PID 1 successfully}

\subsection*{Test 3: Register Child Processes}
\textbf{Input:} Two child processes spawned via \texttt{fork()}; both registered with \texttt{register\_pid()}. \\
\textbf{Expected:} Returns 0 for both, children are added to the monitored list. \\
\textbf{Result:} PASS --- \texttt{Registered PID 3860 successfully}, \texttt{Registered PID 3861 successfully}

\subsection*{Test 4: Fetch with 4 Processes Monitored}
\textbf{Input:} \texttt{fetch\_ctx\_switches(\&stats)} with PIDs 3859, 1, 3860, 3861 registered. \\
\textbf{Expected:} Returns cumulative context switches across all four processes and their threads. \\
\textbf{Result:} PASS

\begin{lstlisting}
Voluntary Context Switches:   65097
Involuntary Context Switches: 2523
\end{lstlisting}

\textbf{Verification:} The high voluntary count is largely contributed by PID 1 (systemd), which frequently sleeps waiting for events. The child processes contribute additional switches from their busy work and \texttt{usleep()} calls.

\subsection*{Test 5: Deregister Child 1}
\textbf{Input:} \texttt{deregister\_pid(3860)} \\
\textbf{Expected:} Returns 0, child is removed from the list. \\
\textbf{Result:} PASS --- \texttt{Deregistered PID 3860 successfully}

\subsection*{Test 6: Fetch with 3 Processes Monitored}
\textbf{Input:} \texttt{fetch\_ctx\_switches(\&stats)} with PIDs 3859, 1, 3861 remaining. \\
\textbf{Expected:} Lower cumulative count after removing one child. \\
\textbf{Result:} PASS

\begin{lstlisting}
Voluntary Context Switches:   65086
Involuntary Context Switches: 2520
\end{lstlisting}

\textbf{Verification:} Both counts decreased compared to Test 4 (voluntary: 65097 $\rightarrow$ 65086, involuntary: 2523 $\rightarrow$ 2520), confirming PID 3860's contribution was removed.

\subsection*{Test 7: Deregister Child 2}
\textbf{Input:} \texttt{deregister\_pid(3861)} \\
\textbf{Expected:} Returns 0, child is removed from the list. \\
\textbf{Result:} PASS --- \texttt{Deregistered PID 3861 successfully}

\subsection*{Test 8: Fetch with 2 Processes Monitored}
\textbf{Input:} \texttt{fetch\_ctx\_switches(\&stats)} with PIDs 3859 and 1 remaining. \\
\textbf{Expected:} Further decrease after removing second child. \\
\textbf{Result:} PASS

\begin{lstlisting}
Voluntary Context Switches:   65075
Involuntary Context Switches: 2519
\end{lstlisting}

\textbf{Verification:} Counts decreased again (voluntary: 65086 $\rightarrow$ 65075, involuntary: 2520 $\rightarrow$ 2519), confirming PID 3861's contribution was removed.

\subsection*{Test 9: Deregister Current Process}
\textbf{Input:} \texttt{deregister\_pid(3859)} \\
\textbf{Expected:} Returns 0, current process is removed from the list. \\
\textbf{Result:} PASS --- \texttt{Deregistered PID 3859 successfully}

\subsection*{Test 10: Fetch with Only PID 1}
\textbf{Input:} \texttt{fetch\_ctx\_switches(\&stats)} with only PID 1 remaining. \\
\textbf{Expected:} Returns context switches for PID 1 alone. \\
\textbf{Result:} PASS

\begin{lstlisting}
Voluntary Context Switches:   65074
Involuntary Context Switches: 2519
\end{lstlisting}

\textbf{Verification:} Voluntary count dropped by 1 (65075 $\rightarrow$ 65074), confirming the test process was properly removed. The clear downward trend across Tests 4, 6, 8, and 10 validates that \texttt{deregister} correctly removes processes from the monitored list.

\subsection*{Test 11: Deregister PID 1}
\textbf{Input:} \texttt{deregister\_pid(1)} \\
\textbf{Expected:} Returns 0, PID 1 is removed. \\
\textbf{Result:} PASS --- \texttt{Deregistered PID 1 successfully}

\subsection*{Test 12: Invalid PID}
\textbf{Input:} \texttt{register\_pid(-1)} \\
\textbf{Expected:} Returns $-1$ with \texttt{errno} set to \texttt{EINVAL}. \\
\textbf{Result:} PASS --- \texttt{register\_pid failed: Invalid argument}

\subsection*{Test 13: Deregister Non-Existent PID}
\textbf{Input:} \texttt{deregister\_pid(99999)} \\
\textbf{Expected:} Returns $-1$ with \texttt{errno} set to \texttt{ESRCH}. \\
\textbf{Result:} PASS --- \texttt{deregister failed: No such process}

\subsection*{Test 14: Double Deregister}
\textbf{Input:} \texttt{deregister\_pid(3859)} (already deregistered in Test 9). \\
\textbf{Expected:} Returns $-1$ with \texttt{errno} set to \texttt{ESRCH}. \\
\textbf{Result:} PASS --- \texttt{deregister failed: No such process}

\textbf{Verification:} Confirms that a PID cannot be deregistered twice, and the list correctly reflects prior removals.

\section*{Summary}

\begin{table}[h]
\centering
\begin{tabular}{@{}llc@{}}
\toprule
\textbf{Task} & \textbf{Test Case} & \textbf{Result} \\
\midrule
A & Process-specific mode (current process) & PASS \\
A & Process-specific mode (PID 1) & PASS \\
A & System-wide mode & PASS \\
A & Invalid PID $\rightarrow$ ESRCH & PASS \\
A & Invalid flags $\rightarrow$ EINVAL & PASS \\
A & Buffer too small $\rightarrow$ ENOSPC & PASS \\
\midrule
B & Register current process & PASS \\
B & Register PID 1 & PASS \\
B & Register child processes & PASS \\
B & Fetch (4 processes monitored) & PASS \\
B & Deregister child 1, fetch (3 monitored) & PASS \\
B & Deregister child 2, fetch (2 monitored) & PASS \\
B & Deregister self, fetch (1 monitored) & PASS \\
B & Deregister PID 1 & PASS \\
B & Invalid PID $\rightarrow$ EINVAL & PASS \\
B & Deregister non-existent PID $\rightarrow$ ESRCH & PASS \\
B & Double deregister $\rightarrow$ ESRCH & PASS \\
\bottomrule
\end{tabular}
\end{table}

\noindent\textbf{All 17 tests passed.}

\end{document}